\documentclass[11pt]{article}

\usepackage[review]{acl}
\usepackage{times}
\usepackage{latexsym}
\usepackage[T1]{fontenc}
\usepackage[utf8]{inputenc}
\usepackage{microtype}
% inconsolata is optional in the ACL template and is not required to build.
\usepackage{amsmath,amssymb,booktabs}

\title{DECORD: Repairing Option-Letter Commitment in Multiple-Choice Reasoning}

\author{Anonymous}

\begin{document}
\maketitle

\begin{abstract}
Small instruction-tuned models often calculate a quantity that matches one option and then emit a different letter.
We separate that failure from not knowing the answer, and we study two repairs that share the same decision rule.
DECORD is a training-free commitment step: the emitted letter is kept unless the trace already names an option, in which case the next-token distribution at a forced answer boundary is shifted toward that option.
DECORD-FT trains a LoRA adapter with extra weight on the answer line, a margin between the gold letter and the other letters, and a reference-model preference loss on the model's own wrong endings.
On AQUA-RAT, both methods are compared with the untuned model, uniform supervised fine-tuning, and standard DPO under a dev-selected hyperparameter freeze.
The registered test comparison is reported in Table~\ref{tab:main}.
\end{abstract}

\section{Introduction}

Multiple-choice algebra items punish two different mistakes with the same zero.
The model may never reach the right quantity, or it may reach that quantity and still name the wrong letter.
The second mistake is the one this paper targets.
It is easy to miss, because a letter extracted from a long trace can be an intermediate mention of an option rather than a commitment, and because a fine-tune that raises the chance of the correct rationale can still leave the final letter behind.

We treat the final letter as its own decision.
DECORD reads a finished trace, checks whether the numbers written before the answer line match exactly one option, and changes the letter only when they do.
DECORD-FT teaches the same commitment during LoRA training, so the greedy letter moves toward the option the rationale already supports.
The comparison is deliberately narrow: the same questions, the same prompt, the same LoRA budget, and preference pairs built once from the supervised model and reused by both DPO and DECORD-FT.

Inference-time steering of a contrastive direction can change what a model says without retraining it \citep{li2023iti}.
Option order and option wording also move multiple-choice accuracy even when the underlying fact does not \citep{zhao2021calibrate,zheng2024mcq}.
DECORD is closer to the second line of work than to a general truthfulness direction.
The signal it uses is the model's own arithmetic, not a probe trained to detect honesty.
The training loss is standard preference optimization \citep{rafailov2023dpo} plus an explicit letter margin, not a reward model.

\section{Method}

\subsection{Task}

Each item is a question with options \(A\)--\(E\) and a gold letter \(y\).
The model writes a trace under a fixed chat prompt and we read a letter \(\hat{y}\).
Accuracy is \(\hat{y}=y\).
A trace has a \emph{numeric commitment} when a number in the pre-answer span matches exactly one option's value; call that letter \(z\).
A mismatch is the event \(z\neq\hat{y}\) when \(z\) exists.
Items with no numeric commitment are ordinary errors, not mismatches.

\subsection{DECORD}

Given one trace, let \(\ell(L)\) be the log-probability of letter \(L\) as the next token after the prefix that ends in \texttt{Final answer:}.
If \(z\) is empty, DECORD returns the letter already emitted in the trace, or \(\arg\max_L \ell(L)\) when the trace has no letter.
If \(z\) is present,
\[
s(L)=\ell(L)+\alpha\cdot\begin{cases}+1 & L=z\\ -1 & L\neq z\end{cases},
\qquad
\hat{y}=\arg\max_L s(L).
\]
\(\alpha=0\) ignores the numeric commitment.
We also draw \(K-1\) extra samples at temperature \(0.7\) and take a majority vote over the \(K\) committed letters.
The pair \((\alpha,K)\) is chosen on a 60-item development slice and then frozen.
\(K=1\) is the single greedy trace.
Ties in the vote break toward the earlier letter in alphabetical order.

\subsection{DECORD-FT}

All trained systems use the same LoRA on Qwen2.5-0.5B-Instruct \citep{qwen2025}: rank 16, \(\alpha_{\mathrm{LoRA}}=32\), dropout 0.05, applied to every attention and MLP projection.
The supervised target is the dataset rationale followed by \texttt{Final answer: } \(y\) \citep{ling2017aqua}.
Uniform SFT minimizes the mean token loss on the completion.
DECORD-SFT uses the same data and the same number of steps, multiplies the loss on the last four completion tokens by 4, and adds a hinge of weight \(0.5\) that keeps the gold letter logit at least one unit above each other letter logit at the answer position.

Rejected traces are generated from the uniform SFT model.
If that sample already has the gold letter, its ending is replaced by a wrong letter so the pair still contrasts a correct ending with an incorrect one.
Standard DPO starts from uniform SFT.
The reference log-probabilities are the SFT model's sequence log-probabilities of the chosen and rejected completions, cached before any DPO step, and the loss is
\[
-\log\sigma\Big(\beta\big[(\log\pi_\theta-\log\pi_{\mathrm{ref}})_w-(\log\pi_\theta-\log\pi_{\mathrm{ref}})_l\big]\Big)
\]
with \(\beta=0.1\).
DECORD-DPO starts from DECORD-SFT, uses the same pairs and the same reference procedure on its own initial weights, and adds the letter hinge on the chosen sequence with weight \(0.25\).
Development accuracy chooses whether the fine-tuned system is DECORD-SFT or DECORD-DPO.
That choice is frozen before the test split is scored.

\section{Experimental setup}

The model is Qwen2.5-0.5B-Instruct in float16 on one RTX 2080 Ti.
Decoding for the main table is greedy with 256 new tokens.
Training uses sequences truncated to 512 tokens from the left of the prompt and the right of the completion so the final answer line is kept, one epoch, batch size 1, and gradient accumulation 8.
The training pool is a 400-item sample of AQUA-RAT train, seed 42.
Hyperparameters are selected on 60 development items that do not overlap the training pool.
The test set is the full AQUA-RAT test split (254 items) and is not used to choose \(\alpha\), \(K\), or the fine-tuned variant.

The baselines are the untuned model, uniform SFT, and standard DPO.
DECORD is the training-free rule on the untuned model.
The fine-tuned scheme is the dev-selected DECORD adapter, scored with greedy decoding.
A result counts only if both DECORD and the fine-tuned scheme exceed the best of the three baselines by at least one percentage point on test accuracy.

\section{Results}
\label{sec:results}

Table~\ref{tab:main} will be filled from the registered test run.
Until that file is frozen, this draft does not contain accuracy numbers.

\begin{table}[t]
\centering
\small
\begin{tabular}{lcc}
\toprule
Method & Accuracy & Mismatch \\
\midrule
Base & -- & -- \\
SFT & -- & -- \\
DPO & -- & -- \\
DECORD & -- & -- \\
DECORD-FT & -- & -- \\
\bottomrule
\end{tabular}
\caption{AQUA-RAT test letter accuracy and mismatch rate. DECORD uses the development-selected \((\alpha,K)\). DECORD-FT is the development-selected adapter.}
\label{tab:main}
\end{table}

\section{Limitations}

The comparison uses one model size and one algebra multiple-choice set.
A one-point gain on 254 items is a handful of questions, so the table has to be read with that width in mind.
Numeric commitment misses equivalent forms that do not match as floats, and a wrong intermediate number can pull the letter to a distractor.
Four samples do not exhaust test-time scaling.
The preference pairs reuse one SFT sample per question rather than a searched set of rejections.
Training-free DECORD cannot recover an item whose trace never states the right quantity.

\bibliography{custom}

\end{document}
